\chapter{Glossary (Plain-Language)}
\label{app:glossary}

This glossary explains the recurring technical terms of the dissertation
in everyday language. Definitions are intentionally informal; the precise
versions live in Chapter~\ref{ch:methodology}.

\begin{description}

  \item[Agent] The learning program --- the ``player.'' Here: a small
  neural network deciding each day whether to hold, buy, or sell.

  \item[Environment] The world the agent acts in. Here: a simulated
  trading month built from real historical prices, which executes the
  agent's trades, charges fees, and keeps the books.

  \item[State] Everything the agent is shown before deciding --- its
  ``eyes.'' Here: nine numbers covering recent market behaviour, how far
  through the month it is, and its own cash and holdings.

  \item[Action] One of the choices available. Here there are exactly
  three: do nothing, buy a fixed money slice of the asset, or sell one.

  \item[Reward] The score received after each action. Here: the day's
  change in total wealth, with trading fees already deducted. The agent's
  whole purpose is to make the month's total score as large as possible.

  \item[Episode] One complete round of the task, after which everything
  resets. Here: one calendar month of trading, ending with a forced sale
  of any remaining holdings.

  \item[Policy] The agent's decision rule --- its ``habits.'' Formally a
  function from state to action probabilities; informally, what it would
  do in any situation.

  \item[Policy gradient / REINFORCE] A family of methods that improve
  the policy directly: after each month, decisions that preceded a good
  outcome are made more probable, and vice versa. REINFORCE is the
  original, simplest member.

  \item[Q-value / Deep Q-learning] The other family: learn a ``price
  list'' saying how much future profit each action is worth in each
  situation, then take the best-priced action. ``Deep'' means the price
  list is a neural network rather than a lookup table.

  \item[PPO] A modern, stabilised descendant of REINFORCE that limits how
  much the decision rule can change per update --- steadier, but capable
  of settling into whichever behaviour it found first.

  \item[Neural network] A flexible mathematical function with many
  adjustable dials (parameters), tuned gradually so its outputs improve
  against an objective.

  \item[Training / validation / test split] The data is cut into three parts
  by date. The agent learns on the earliest (2018--2022); every tuning
  decision is made by looking only at the middle part (2023); and the final
  examination uses the latest part (2024--2025), which nothing in the design
  was allowed to see. The middle part is what makes the last part a real
  exam rather than a rehearsal: without it, the choices would have been made
  by peeking at the answers.

  \item[Baseline] A simple non-learning strategy every agent must be
  compared against. The important one here is \emph{buy-and-hold}:
  invest everything at the start of the month and wait.

  \item[Buy-and-hold] The classic passive strategy. Decades of evidence
  say it is genuinely hard to beat using price history alone --- which is
  why beating it easily in an experiment is grounds for suspicion, not
  celebration.

  \item[Transaction fee / basis points (bps)] The cost of trading,
  charged as a fraction of the traded amount. 5 bps $=$ 0.05\%, i.e.\
  50 cents per \$1{,}000 traded.

  \item[Sharpe ratio] Profit earned per unit of ups-and-downs endured.
  Higher is better, but a strategy that barely invests can post a
  flattering Sharpe simply by barely moving.

  \item[Maximum drawdown] The worst peak-to-trough fall an investor
  following the strategy would have suffered --- a measure of the most
  painful stretch, not the average one.

  \item[Exposure] How much of the capital is invested in the asset at a
  given moment, from 0\% (all cash) to 100\% (fully invested).

  \item[Seed] The starting value of the random number generator. Running
  the same experiment with several seeds checks that a result is not a
  fluke of one lucky roll.

  \item[Overfitting] Learning the quirks of the practice data rather
  than the lesson, and therefore failing on new data. The chronological
  split and held-out months guard against it.

  \item[Masking] Greying-out impossible actions (selling what you do not
  own; buying with no cash) so the agent can neither attempt nor learn
  from illegal moves.

  \item[Mask-determined (state-blind) policy] An agent whose choices can be
  predicted from the greyed-out actions alone, without looking at anything
  else it was shown. Such an agent is following a habit --- ``buy whenever
  buying is allowed'' --- rather than making a decision, and in a rising
  market the habit scores well, which is why profit alone cannot reveal it.
  Detecting this is what Section~\ref{sec:res_conditioning} is for.

  \item[Admissible state] A state description that is legally usable in this
  framework: the score can be worked out from it, the next state follows from
  it, and it says enough to choose an action. The test has teeth --- it is the
  reason drawdown had to be added to the state before it could be added to
  the reward.

  \item[Accounting identity] The bookkeeping check that the sum of all
  daily scores in a month exactly equals the month's true profit. It is
  verified in every single experiment; if it ever failed, the simulator
  would be lying to the agent.

  \item[Markov property] The requirement that the state contain
  everything relevant from the past --- like a chess photograph being
  enough to resume the game. Much of this project's design effort went
  into honouring it.

  \item[Partial observability (POMDP)] The honest admission that no
  price chart contains everything that moves markets (news, order flow,
  other traders' intentions). The theory for acting under such hidden
  information exists but is much harder; here it is acknowledged as a
  limitation and a future direction.

\end{description}
